\section{Discussion and Limitations}
\label{sec:discussion}

\subsection{What the Method Delivers}

The canonical method uses a single pretreatment proxy in two roles that are available from observed data: supervising an outcome-free representation and measuring the residual imbalance of $(X,W)$ across treatment arms within that representation. Under uniform outcome-relevant stability, the product $\Gamma D(\phi)$ is a sensitivity certificate for hidden-confounding bias. This is stronger than reporting a conditional-independence test alone because it targets a causal error scale, but weaker than point identification when $D(\phi)>0$ because $\Gamma$ is not generally identified from one proxy.

\subsection{Testable and Structural Components}

Held-out joint balance, empirical overlap, and minimum stratum mass are observable properties. Proxy-channel invariance, conditional mixture injectivity, and outcome-relevant stability are structural assumptions about the full-data law. Exact balance is therefore a testable endpoint only up to sampling error, while its causal interpretation still requires structural assumptions. A large held-out discrepancy can falsify exact balance; a small test statistic cannot prove exact balance, injectivity, or stability.

\subsection{Attainability and Realized Bias}

For $Z=\phi(X)$ with finitely many strata, exact balance is attainable only on a narrow boundary characterized in Section~\ref{sec:identification}. This is a strength as a safety certificate and a limitation as a universal identification method. In the approximate regime, reducing $D(\phi)$ reduces the certified worst-case bias under a common $\Gamma$, but it need not reduce the realized bias without additional alignment conditions. Cancellation, outcome-irrelevant proxy directions, bias floors, and ill-posed channels explain why the two orderings can differ.

\subsection{Learning and Inference Gaps}

The current finite-sample theorem covers enumerable representation classes and fixed-fold evaluation. The neural experiments use soft assignments and covariance-weighted surrogate objectives that are not algebraically identical to the canonical MMD aggregation. A neural-class generalization bound, a soft-to-hard guarantee, and valid joint inference after adaptive representation selection remain open. These gaps prevent the manuscript from describing every neural result as a direct verification of the theory, even when its empirical error is lower.

\subsection{Audit Views and Practical Scope}

Cross-mask learning uses one proxy coordinate for audit and the remaining high-dimensional measurement for representation learning, then averages across audit orientations. This design blocks literal self-copying of the audited coordinate. It does not by itself establish conditional independence among proxy views, so point identification still needs a multi-view structural assumption and an audit-channel null-space condition. The WSC experiment further shows that an encouraging benchmark gap can coexist with a failed practical-overlap rule. Real-data evidence must therefore report overlap, audit discrepancies, and benchmark uncertainty without treating any one diagnostic as ground truth.

\subsection{Scope of the Contribution}

The defensible contribution is the combination of a proxy-supervised balancing representation, an observable discrepancy, an exact certification boundary, and a stability-indexed approximate guarantee. The paper does not claim that one generic proxy point-identifies the ATE, that every smaller discrepancy has smaller realized bias, or that the current neural optimizer reaches the population optimum. A systematic literature audit and broader external validation are required before any firstness or universality claim.
